\begin{resumo}
 O projeto v-FES cycling, realizado pela Plataforma Tecnológica BEM-TE-VI da Universidade de Brasília, acompanha
 seus participantes por meio de avaliações regulares feitas por uma equipe de fisioterapeutas que utiliza sete
 instrumentos. O aplicativo Bem-Te-Vi reúne esses formulários, mas trabalha com cada um deles separadamente, só
 oferece histórico e comparação das avaliações no fluxo da CIF e apresenta problemas no registro de dados. Este
 trabalho define os requisitos e um modelo de dados interoperável em HL7 FHIR para estruturar esses dados e
 apoiar o acompanhamento das avaliações, tendo como ponto de partida o módulo CIF implementado pelo autor. A
 pesquisa segue o método de \textit{design science} proposto por Wieringa e baseia-se no
 código do aplicativo, na documentação do projeto e na literatura especializada sobre esses instrumentos, sem
 consultas à equipe de fisioterapeutas. Foram especificados 67 requisitos, e o modelo, criado em FHIR R4,
 representou todas as estruturas da CIF, com instâncias de dados fictícios aprovadas pelo validador. A análise
 mostrou ainda que o aplicativo já automatiza os cálculos em cada avaliação, embora parte deles se desvie do
 instrumento, mas não automatiza o acompanhamento entre as avaliações nos demais formulários. Esses resultados
 servem como orientação para a implementação do banco de dados e a validação com a equipe de fisioterapeutas no
 TCC 2.

 \vspace{\onelineskip}

 \noindent
 \textbf{Palavras-chave}: Engenharia de requisitos. Dados clínicos. CIF. HL7 FHIR. Fisioterapia.
\end{resumo}
